\documentclass[11pt]{article}
\usepackage[top=1.00in, bottom=1.0in, left=1.1in, right=1.1in]{geometry}
\usepackage{graphicx}
\usepackage{natbib}
\usepackage{amsmath}
\usepackage{lineno}
\usepackage{xr-hyper}
\externaldocument{..//chillingms}
\newcommand{\R}[1]{\label{#1}\linelabel{#1}}
\newcommand{\lr}[1]{line~\lineref{#1}}
\setlength\parindent{0pt}

\parskip=5pt

\begin{document}
Editor and reviewer comments (we provide below the full context of each review) are in \emph{italics}, while our responses are in regular text. \\ 

{\bf Editor's comments:} 

% NPH-V-2026-58419 - Viewpoint -  The problem and promise of modeling plant leafout under warmer winters
% by Wolkovich, Elizabeth; Ngo, Justin; Van der Meersch, Victor; Guy, Robert; Auerbach, Jonathan

\emph{I am sorry for the long time needed to review your manuscript. We have been waiting for third review, which was not delivered. Should it be delivered, we will forward it to you as well and ask you to consider it for a revision.}

\emph{Looking at the reports, appended at the end of this message, you can see quite a number of critical issues that
need to be addressed. Therefore, as it stands, your paper is not suitable for publication. However, I believe
that a careful but massive revision can solve the problems indicated by the reviewers (esp. ref. 2).}

\emph{Please revise your viewpoint and return your paper within the next 8 weeks. It will be send out for a second
round of reviewing to support a final decision.}
% Sincerely,
% Andrea Polle
% Editor, New Phytologist

{\bf Referee comments}

{\bf Referee 1} 

\emph{I found this Viewpoint very interesting and valuable for readers interested in either modelling or phenology. The
manuscript provides a useful overview of the field and important recommendations. However, there are several
areas where clarification, additional explanation, and revision would strengthen the work, particularly in the
figures and figure legends.}

\emph{Specific Comments}

\emph{•       A clearer distinction between observational and experimental data should be introduced earlier in the
manuscript, perhaps within the Introduction. At present, Figure 1 appears to refer primarily to experimental
data, whereas Figure 2 appears to use observational data, but this distinction is not clearly established when
the figures are first introduced. Clarifying this earlier would reduce potential confusion for the reader.}

Good point, we reworked the introduction and now introduce both data types in the second paragraph, see edits \R{r1expobsdatS}-\R{r1expobsdatE}.

\emph{•       Line 29: The statement regarding "choice of parameterization" yielding different results requires further
clarification. Such outcomes are not uncommon in modelling studies, so it would be helpful to explain more specifically what is being referred to here. For example, are the parameters inadequately constrained?}

This is a good point and we have now modified the sentence to give two examples (see \lr{r1clarify}), since many issues can lead to this. 

\emph{Figure 2:}

\emph{  •	The legend is unclear, particularly the phrase "assuming ... is observed". Additional explanation is needed.}

\emph{  •	It would be helpful to specify the temperature profile used in the example.}

\emph{  •	Figure 2 presents estimation from a single winter. It would be useful to explain how the conclusion about indeterminacy would differ if the estimation was made using multiple conditions, as those illustrated in Figure 1. In addition, some discussion of how frequently models are estimated from a single condition versus multiple conditions would provide useful context.}

\emph{•	Line 178: Arabidopsis thaliana is referred to as "biennial" but is normally an annual.}

\emph{•	Line 179, Wilczek et al 2009, Kim et al 2009: Would a more recent review of Arabidopsis vernalization be appropriate here? }

\emph{•	Figure 3: Figure 3 requires substantial revision. At present, it is not explained in sufficient detail and presents the florigen pathway in a potentially misleading manner. In particular:}
\emph{ •	FLC is a flowering repressor,}
\emph{ •	COOLAIR is cold-induced,}
\emph{  •	FT is absent from the figure, despite being the canonical florigen gene.}

\emph{  •	CONSTANS protein is regulated by photoperiod and is the only circadian clock-related gene included in panel 4.}

\emph{  •	Additional explanation is needed for the genes shown in the other panels. Readers may find it difficult to understand their relevance without further context.}

\emph{  •	It would also be helpful to state explicitly which species each pathway or gene network applies to.}

\emph{  •	More broadly, the authors may wish to consider whether all of the detail included in the figure is necessary for the aims of the article.}

\emph{  •	Panel 1 is not entirely clear in its current form. I suggest splitting the panel to better reflect the two types of models discussed, namely sequential versus parallel models.}

\emph{•	Line 261: The distinction between the different meanings of the term "chilling" is unclear. The statement regarding uncertainty about the physiological phase requires further explanation. It would be useful to clarify whether this reflects a limitation of the methods or a feature of the experimental design.}

\emph{Figure 4:}

\emph{  •	Figure 4 would benefit from a more detailed explanation. At present, the interpretation of the figure is not sufficiently clear.}
  
\emph{  •	I could not find a definition of GDD (Growing Degree Days) within the manuscript, although the term appears in the figure. The abbreviation should be defined when first introduced.}
  
\emph{•	Supplementary Material: The supplementary material describes "quantitative PCR (one type of transcriptomics)". This statement is not correct. Transcriptomics refers to approaches that assay the complete set of transcripts, whereas qPCR is a targeted technique used to quantify the expression of specific genes of interest. The wording should therefore be revised.}


{\bf Referee 2} 

\emph{This viewpoint provides an interesting read and raises a relevant and well-known problem in chilling-model literature : chilling is not directly observable and modelling chilling requirements requires assumptions about when chilling accumulation begins as well as the temperature range within which chilling is considered effective. In many models, these assumptions involve defining a starting date and minimum and maximum temperature threshold for chill accumulation. I don’t quite get why the starting date is such an issue for the authors as from my experience chill starts when temperature fall in the range within which chilling is considered effective whether one starts counting in September, October of November, but maybe this point is more important in higher latitude ?}

\emph{I have two main concerns regarding this paper.}

\emph{My first concern should be relatively straightforward to address. In my opinion, several key references from the literature are missing, while a number of references from less established sources are cited (ex : MDPI journals). I provide below some references that I believe are particularly relevant to the topic and encourage the authors to review the cited literature more broadly and ensure that the key contributions to the field are adequately represented.}

\emph{My second concern is that some key information is missing, which makes parts of the viewpoint difficult to understand for a broader audience. For example, the diversity of chilling models, how they calculate chilling (chilling hours, chilling units, and chilling portions) and their different assumptions regarding chilling accumulation, including potential reversal or negation of previously accumulated chilling, is not sufficiently described. I believe that Figure 1 or Figure 2 could be substantially revised to provide a more informative overview of these different modelling approaches and their underlying assumptions. Figure 1 is not so informative in its current state. Following the same line of thought, the authors discuss at length the limitations of process-based models and emphasize the need to provide stronger mechanistic foundations for these models. However, they do not sufficiently discuss which biological processes could provide such mechanistic foundations and don’t mention that some models include mechanistic processes. The authors mention callose and hormones as key molecular basis of chilling, but the underlying process could be better described. Similarly, other important candidates, such as DAM genes, are not discussed in the main text. In addition, Figure 3 is not, in my opinion, used to its full potential: the genes presented in the figure are not discussed in the text, which makes it difficult to understand their relevance to the argument being developed, especially for non-experts. The figure legend could also be revised to make the figure more accessible and informative. I noticed that the supplementary information does provide a more detailed description of these molecular indicators. I believe that a middle ground could be found by bringing some of this information into the main text, while retaining the more detailed discussion in the Supplementary Information. Finally, while callose deposition can contribute to the temporary closure of plasmodesmata and thereby restrict symplastic communication, I would be interested in hearing the authors' perspective on which specific biological processes may be regulated by hormones and key transcription factors during dormancy progression and release. In particular, it would be valuable to discuss how these molecular regulators could be linked to the physiological processes that are ultimately represented in chilling and dormancy models. I understand these viewpoint should be forward thinking, and if the authors point is that “chilling” is obsolete, it is legitimate to make a clear proposal of what should replace it.}

\emph{I provide some minor comments below that I hope will help the authors address these two main concerns.}

\emph{L38 : I think it useful in this sentence to clarify that spring phenology is important for crops and associated food safety/food production.}

We have added this point here, see \R{r2foodsafety}.

\emph{L46 : I appreciate the perspective presented in this section. However I think it would be useful to provide a brief description of how these model work, including the concepts of chilling hours, chilling unit, chilling portion etc ,so that non-expert readers can better understand the different approaches and form their own opinion based on the information provided. Throughout the text, chilling units are mentioned, but no definition is provided.} 

\emph{L54 : I don’t really understand what the authors mean with “a life history event”.}

\emph{L58-61 : sentence is slightly too informal. I advise to rephrase. Also there is some literature on the shared mechanisms of seed dormancy and bud dormancy whereas from this sentence it may seem that shared mechanisms are just likely. Clarify, read and include useful literature , maybe : }

\emph{Leida et al. (2012), PLOS ONE — ``Chilling-Dependent Release of Seed and Bud Dormancy in Peach Associates to Common Changes in Gene Expression."}

\emph{Tuan PA, Bai S, Saito T, Ito A, Moriguchi T. Dormancy-Associated MADS-Box (DAM) and the Abscisic Acid Pathway Regulate Pear Endodormancy Through a Feedback Mechanism. Plant Cell Physiol. 2017 Aug 1;58(8):1378-1390. doi: 10.1093/pcp/pcx074. PMID: 28586469.}

\emph{Sato \& Yamane (2024), Journal of Experimental Botany — ``Histone modifications affecting plant dormancy and dormancy release: common regulatory effects on hormone metabolism."}

\emph{L65-68 : I find this sentence a bit misleading as one may understand that only some parts of the cuttings are exposed to temperature treatment. Further, classic Weinberger approach are temperature  0-7,2 °C, not trying different temperature while here one can understand different temperatures are tested. Finally, Ettinger might not be the best reference here as it is a metanalysis and therefore not the best citation for the cuttings experiment. Maybe Vitasse et al., 2014 might be a better one ?}

\emph{L72-76: Sentence is long and hard to follow, remove “ and has led to the cross-disciplinary field of phenology.}

\emph{L78-83 : This section is not particularly insightful in its current form and i suggest to shorten, or remove. In addition, the term “molecular approach” is sometimes used by the authors, although the experimental work described in these sentences does not necessarily fall within the scope of a molecular approach. In particular, using Arabidopsis thaliana as an experimental system does not, in itself, imply that the work is molecular biology. I would therefore suggest avoiding such broad generalizations and using terminology that more accurately reflects the approaches described. More generally, genetics and molecular biology encompass distinct, although overlapping, fields, and I think it would be useful to distinguish between them more precisely throughout the text.}

\emph{L89-94 : I think here would be a good place to have clarifications on models, what chilling units, hours and portion mean, etc. }

\emph{140 : Could you please clarify this statement? In particular, it would be helpful to explain why the data are considered uninformative, and what aspects of the data collection or quality might be improved to make them more useful for the analysis. More broadly, to what extent can limitations in the data explain the models’ predictive performance ? Could some of the observed limitations instead be related to the model specification or methodology ?}

\emph{147 : I really like this section.}

\emph{L191-195 : More mechanistic models based on enzymatic response to temperature have been proposed and should be cited. I provide example below, and let the authors select and justify what they think is relevant. Also more information on figure 3 and molecular basis, associated physiological process would be useful here.}

\emph{Luedeling et al. 2021  “Reducing the uncertainty on chilling requirements for endodormancy breaking of temperate fruits by data-based parameter estimation of the dynamic model: a test case in apricot.” It might be a good source because this paper explicitly explains that the Dynamic Model assumes that endodormancy is completed when a certain amount of a dormancy-breaking factor (DBF) has accumulated.}

Also :
\emph{Huanjiong Wang, Wenrui Bai, Zhi Hu, Shaozhi Lin, Quansheng Ge, Effectiveness of freezing temperatures on dormancy release of temperate woody species, Annals of Botany, Volume 134, Issue 4, 1 October 2024, Pages 615–630, https://doi.org/10.1093/aob/mcae112}

\emph{Sperling, O., Kamai, T., Tixier, A., Davidson, A., Jarvis-Shean, K., Raveh, E., DeJong, T., \& Zwieniecki, M. A. (2019).``Predicting bloom dates by temperature mediated kinetics of carbohydrate metabolism in deciduous trees."}

\emph{Charrier, G., Lacointe, A. \& Améglio, T. (2018). Dynamic Modeling of Carbon Metabolism During the Dormant Period Accurately Predicts the Changes in Frost Hardiness in Walnut Trees (Juglans regia L.).Frontiers in Plant Science 9:1746.}

\emph{220-223 : I don’t understand this sentence, and don’t understand the link with Figure 4. Unfortunately, I can’t help because I don’t understand their point . Revise and clarify this point}

\emph{233 : avoid `molecular biologists', especially for work that do not include molecular biology.}

\emph{Figures :}

\emph{Figure 1 : I disagree with this definition of chilling : « The term chilling generally refers to a process where dormant embryonic tissues exposed to cool temperatures accelerate a plant event that later occurs » . Insufficient Chilling generates abnormal flower and can have detrimental effect on yield ( see example   Enrique González-Martínez, Jaume Lordan, Sergi Munné-Bosch, Xavier Miarnau,Impact of rising temperatures during dormancy on key yield components in almond,Scientia Horticulturae,Volume 355,2026,). I don’t think chilling only accelerates process, it contributes to the proper resumption of bud development and has importance for the reproductive success of species. This figure could be improved by clarifying that chilling duration is modulated by experimentation, especially since there is no y axis. Further, while I understand that Tmax and Tmin can fluctuate, providing range of temperatures for the chilling and forcing period would make the figure clearer for a broader audience and easier for this audience to grasp the concepts.}

\emph{Figure 3 is mentioned prior figure 2 , also it would be better to have consistency with Fig. 1or Figure 1.}

\emph{Figure 2 : ``In both cases, multiple parameter sets reproduce the same observation equally well."  I would remove this sentence.}

\emph{Figure 3 legend needs to be extended. }

\emph{``Over 30 variants of this model exist, including those where the phases are sequential (darker shading) or occur in parallel (dark and light shading).'' Consider placing some models on this figure so one can understand this diversity of models. Also there is a disconnect between the main text and figure content. I think some of its contents needs to be referred to in the text, especially the potential molecular basis of chilling. I would clarify `gene expression' in the y-axis and provide title for molecular process : Callose deposition, GA , Circadian clock and flowering time. ``Molecular work suggests callose (second panel), alongside gibberellic acid (GA, third panel) and circadian clock-related genes (fourth panel) may all underlie these hypothesized phases.'' i think one needs to clarify if this is hypothetical expression trend or based on observed trend ( then studies should be cited).}

%\bibliography{..//..//..//refs/decsens.bib}
%\bibliographystyle{..//..//..//refs/bibstyles/amnat.bst}

\end{document}
